% GENERATED FILE. Edit canonical lessons, then run npm run generate:books.
% canonical-source-hash: fnv1a64:88cfba01669882cc
% canonical-lessons: LA-C470-ascia, LA-C470-terebra, LA-C470-runcina, LA-C470-lima, LA-C470-dolo, LA-R470-review-463-470

\chapter{Adze, Drill, Plane, File, To hew}
\label{ch:la-adze-drill-plane-file-to-hew}
\hlchaptermodality{470}

\begin{chapteropening}
\textbf{By the end of this chapter:} \emph{I can name a carpenter's tools: an adze, a drill, a plane and a file, and talk about hewing wood.}

Complete the last lesson of chapter 470: I can name a carpenter's tools: an adze, a drill, a plane and a file, and talk about hewing wood.
\end{chapteropening}

\section[ascia, asciae]{ascia, asciae — an adze, a carpenter's axe}

\label{lesson:LA-C470-ascia}



\begin{quote}
Before the new one: say the Latin for a blacksmith, then the Latin for to hammer.
\end{quote}

\subsection*{You'll want to know: ascia, asciae}
\textbf{ascia, asciae} — \textquotedblleft{}an adze, a carpenter's axe\textquotedblright{}.

\begin{quote}
Faber asciam in manū tenet.
\end{quote}

\begin{quote}
\emph{The carpenter holds an adze in his hand.}
\end{quote}

First declension, feminine: \emph{ascia, asciae f.} A tool with the blade across the handle, for shaping beams.

\begin{cousinweb}
\emph{ascia} is inherited; Lewis \& Short compare Greek \emph{axinē} 'axe' as a cousin. EtymDB gives no parent, so the link is hedged.
\end{cousinweb}

\begin{grammarlens}[title={What you've built}]
One more word for telling, explaining and describing at B1.
\end{grammarlens}

\subsection*{Your turn}
\begin{itemize}
\raggedright
  \item \emph{Say it:} \emph{ascia, asciae}
  \item \emph{Say it:} \emph{ascia, asciae}, once more
  \item \emph{Say it:} a sentence of your own with \emph{ascia, asciae}
  \item \emph{Recall:} say the Latin for a blacksmith, then the Latin for to hammer, then say \emph{ascia, asciae} again
\end{itemize}

\begin{tcolorbox}[breakable,colback=teal!4,colframe=teal!35!black,title={Before you move on}]
What does ascia, asciae mean? (\textquotedblleft{}An adze, a carpenter's axe\textquotedblright{}.) Say it once more.
\end{tcolorbox}
\section[terebra, terebrae]{terebra, terebrae — a drill, a borer, an auger}

\label{lesson:LA-C470-terebra}



\begin{quote}
Before the new one: say the Latin for to hammer, then the Latin for an adze.
\end{quote}

\subsection*{You'll want to know: terebra, terebrae}
\textbf{terebra, terebrae} — \textquotedblleft{}a drill, a borer, an auger\textquotedblright{}.

\begin{quote}
Faber terebrā in tabulā labōrat.
\end{quote}

\begin{quote}
\emph{The carpenter works on the board with a drill.}
\end{quote}

First declension, feminine: \emph{terebra, terebrae f.}, from \emph{terō} 'to rub, to wear away'.

\begin{cousinweb}
\emph{terebra} is built on \emph{terō} 'to rub, to wear away' with the tool suffix \emph{-bra} (EtymDB): a tool that bores by turning.
\end{cousinweb}

\begin{grammarlens}[title={What you've built}]
One more word for telling, explaining and describing at B1.
\end{grammarlens}

\subsection*{Your turn}
\begin{itemize}
\raggedright
  \item \emph{Say it:} \emph{terebra, terebrae}
  \item \emph{Say it:} \emph{terebra, terebrae}, once more
  \item \emph{Say it:} a sentence of your own with \emph{terebra, terebrae}
  \item \emph{Recall:} say the Latin for to hammer, then the Latin for an adze, then say \emph{terebra, terebrae} again
\end{itemize}

\begin{tcolorbox}[breakable,colback=teal!4,colframe=teal!35!black,title={Before you move on}]
What does terebra, terebrae mean? (\textquotedblleft{}A drill, a borer, an auger\textquotedblright{}.) Say it once more.
\end{tcolorbox}
\section[runcīna, runcīnae]{runcīna, runcīnae — a plane (a carpenter's tool)}

\label{lesson:LA-C470-runcina}



\begin{quote}
Before the new one: say the Latin for an adze, then the Latin for a drill.
\end{quote}

\subsection*{You'll want to know: runcīna, runcīnae}
\textbf{runcīna, runcīnae} — \textquotedblleft{}a plane (a carpenter's tool)\textquotedblright{}.

\begin{quote}
Faber tabulam runcīnā parat.
\end{quote}

\begin{quote}
\emph{The carpenter prepares the board with a plane.}
\end{quote}

First declension, feminine: \emph{runcīna, runcīnae f.} An iron blade in a wooden block, for smoothing wood.

\begin{cousinweb}
\emph{runcīna} was borrowed from Greek \emph{rhykanē} 'plane', as an ancient grammarian already said (Lewis \& Short); the form was changed to look like \emph{runcō} 'to weed'.
\end{cousinweb}

\begin{grammarlens}[title={What you've built}]
One more word for telling, explaining and describing at B1.
\end{grammarlens}

\subsection*{Your turn}
\begin{itemize}
\raggedright
  \item \emph{Say it:} \emph{runcīna, runcīnae}
  \item \emph{Say it:} \emph{runcīna, runcīnae}, once more
  \item \emph{Say it:} a sentence of your own with \emph{runcīna, runcīnae}
  \item \emph{Recall:} say the Latin for an adze, then the Latin for a drill, then say \emph{runcīna, runcīnae} again
\end{itemize}

\begin{tcolorbox}[breakable,colback=teal!4,colframe=teal!35!black,title={Before you move on}]
What does runcīna, runcīnae mean? (\textquotedblleft{}A plane (a carpenter's tool)\textquotedblright{}.) Say it once more.
\end{tcolorbox}
\section[līma, līmae]{līma, līmae — a file (tool); polishing}

\label{lesson:LA-C470-lima}



\begin{quote}
Before the new one: say the Latin for a drill, then the Latin for a plane.
\end{quote}

\subsection*{You'll want to know: līma, līmae}
\textbf{līma, līmae} — \textquotedblleft{}a file (tool); polishing\textquotedblright{}.

\begin{quote}
Faber ferrum līmā pūrgat.
\end{quote}

\begin{quote}
\emph{The smith cleans the iron with a file.}
\end{quote}

First declension, feminine: \emph{līma, līmae f.} Also the polishing of a poem: \emph{līmae labor}, the work of the file.

\begin{cousinweb}
The origin of \emph{līma} is not settled: Lewis \& Short say plainly that it is unknown, and EtymDB gives no sure parent.
\end{cousinweb}

\begin{grammarlens}[title={What you've built}]
One more word for telling, explaining and describing at B1.
\end{grammarlens}

\subsection*{Your turn}
\begin{itemize}
\raggedright
  \item \emph{Say it:} \emph{līma, līmae}
  \item \emph{Say it:} \emph{līma, līmae}, once more
  \item \emph{Say it:} a sentence of your own with \emph{līma, līmae}
  \item \emph{Recall:} say the Latin for a drill, then the Latin for a plane, then say \emph{līma, līmae} again
\end{itemize}

\begin{tcolorbox}[breakable,colback=teal!4,colframe=teal!35!black,title={Before you move on}]
What does līma, līmae mean? (\textquotedblleft{}A file (tool); polishing\textquotedblright{}.) Say it once more.
\end{tcolorbox}
\section[dolō, dolāre]{dolō, dolāre — to hew, to chop into shape}

\label{lesson:LA-C470-dolo}



\begin{quote}
Before the new one: say the Latin for a plane, then the Latin for a file.
\end{quote}

\subsection*{You'll want to know: dolō, dolāre}
\textbf{dolō, dolāre} — \textquotedblleft{}to hew, to chop into shape\textquotedblright{}.

\begin{quote}
Fabrī trabēs asciīs dolant.
\end{quote}

\begin{quote}
\emph{The carpenters hew the beams with adzes.}
\end{quote}

\emph{dolō, dolāre, dolāvī, dolātum}. Rough shaping of wood with an axe or adze.

\begin{cousinweb}
\emph{dolō} is inherited from Proto-Italic \emph{dolaō} (EtymDB); Lewis \& Short link it with a root meaning 'to split, to tear apart'.
\end{cousinweb}

\begin{grammarlens}[title={What you've built}]
One more word for telling, explaining and describing at B1.
\end{grammarlens}

\subsection*{Your turn}
\begin{itemize}
\raggedright
  \item \emph{Say it:} \emph{dolō, dolāre}
  \item \emph{Say it:} \emph{dolō, dolāre}, once more
  \item \emph{Say it:} a sentence of your own with \emph{dolō, dolāre}
  \item \emph{Recall:} say the Latin for a plane, then the Latin for a file, then say \emph{dolō, dolāre} again
\end{itemize}

\begin{tcolorbox}[breakable,colback=teal!4,colframe=teal!35!black,title={Before you move on}]
What does dolō, dolāre mean? (\textquotedblleft{}To hew, to chop into shape\textquotedblright{}.) Say it once more.
\end{tcolorbox}
\section[(dialogue)]{Words from chapters 463–470 — a review}

\label{lesson:LA-R470-review-463-470}



\begin{quote}
Say the Latin for a vegetable grower and a fishmonger.
\end{quote}

\subsection*{Your turn}
Say these aloud:

\begin{itemize}
\raggedright
  \item a vegetable grower, a fishmonger, a plaited basket, a wicker basket, a bread basket
  \item a steelyard, a pair of scales, a scale pan, to weigh out, to weigh carefully
  \item a peck, a pint, by weight, a quarter-as, a money-changer
  \item the corn supply, dearness, cheapness, to value, to be sold
  \item a cobbler's shop, a barber's shop, an awl, a tanner, a furrier
  \item to dye, the purple shellfish, purple dye, a paint, a paintbrush
  \item an anvil, a pair of bellows, a forge, a blacksmith, to hammer
  \item an adze, a drill, a plane, a file, to hew
\end{itemize}

\begin{tcolorbox}[breakable,colback=teal!4,colframe=teal!35!black,title={Before you move on}]
A vegetable grower? (\textbf{holitor, holitōris}.) A fishmonger? (\textbf{piscārius, piscāriī}.) A plaited basket? (\textbf{sporta, sportae}.) A wicker basket? (\textbf{corbis, corbis}.) A bread basket? (\textbf{canistrum, canistrī}.) A steelyard? (\textbf{statēra, statērae}.) A pair of scales? (\textbf{trutina, trutinae}.) A scale pan? (\textbf{lanx, lancis}.) To weigh out? (\textbf{appendō, appendere}.) To weigh carefully? (\textbf{pēnsō, pēnsāre}.) A peck? (\textbf{modius, modiī}.) A pint? (\textbf{sextārius, sextāriī}.) By weight? (\textbf{pondō}.) A quarter-as? (\textbf{quadrāns, quadrantis}.) A money-changer? (\textbf{nummulārius, nummulāriī}.) The corn supply? (\textbf{annōna, annōnae}.) Dearness? (\textbf{cāritās, cāritātis}.) Cheapness? (\textbf{vīlitās, vīlitātis}.) To value? (\textbf{aestimō, aestimāre}.) To be sold? (\textbf{vēneō, vēnīre}.) A cobbler's shop? (\textbf{sūtrīna, sūtrīnae}.) A barber's shop? (\textbf{tōnstrīna, tōnstrīnae}.) An awl? (\textbf{sūbula, sūbulae}.) A tanner? (\textbf{coriārius, coriāriī}.) A furrier? (\textbf{pelliō, pelliōnis}.) To dye? (\textbf{tingō, tingere}.) The purple shellfish? (\textbf{mūrex, mūricis}.) Purple dye? (\textbf{ostrum, ostrī}.) A paint? (\textbf{pigmentum, pigmentī}.) A paintbrush? (\textbf{pēnicillus, pēnicillī}.) An anvil? (\textbf{incūs, incūdis}.) A pair of bellows? (\textbf{follis, follis}.) A forge? (\textbf{camīnus, camīnī}.) A blacksmith? (\textbf{ferrārius, ferrāriī}.) To hammer? (\textbf{cūdō, cūdere}.) An adze? (\textbf{ascia, asciae}.) A drill? (\textbf{terebra, terebrae}.) A plane? (\textbf{runcīna, runcīnae}.) A file? (\textbf{līma, līmae}.) To hew? (\textbf{dolō, dolāre}.)
\end{tcolorbox}
